\begin{theorem}[任意の有限群値を指定した短いタプル]
\label{thm:group-typed-shortest}
$G=(V,\Sigma,P,S,\fan)$ をfan-out $\le f$，
規則ランク $\le r_{\max}$（$r_{\max}\ge1$）の
reduced good MCFGとする．
$g:\Sigma^*\to A$ を\emph{任意の有限群}
$A$ への準同型とする．各 $B\in V$ と
実現可能な $\mathbf s\in A^{\fan(B)}$ に対し，
\[
 \tuple u\in L_B(G),\qquad
 g^{(\fan(B))}(\tuple u)=\mathbf s
\]
を満たし，かつ
\[
 \|\tuple u\|_1\le
 \beta_g^{\mathrm{grp}}:=
 (r_{\max}+1)\bar\theta_G
 \max\{1,(M_{A,f}-1)|V||A|^f\}
\]
となるタプルが存在する．とくに群と
$f,r_{\max}$ を固定すると
$\beta_g^{\mathrm{grp}}
=O_{A,f,r_{\max}}(\|G\|\bar\theta_G)$ である．
\end{theorem}
\begin{proof}
$B$ から指定した群ベクトルを実現する導出のうち，
（総長，節点数）が辞書式最小のものを選ぶ．
各節点を非終端記号と各成分の群値で標識する．
祖先と真の子孫で同じ標識を持つことはできない．
祖先の部分木を子孫の部分木に取り換えれば，
規則テンプレートを順序どおり評価することで
根までの群値は変わらず，非削除性から総長は増えず，
節点数は真に減るからである．
よってすべての根から葉までの道の節点数は
\(n_g=|V||A|^f\)
以下である．

総出力長が $\bar\theta_G$ を超える節点を重いと呼ぶ．
重い節点は祖先について上に閉じている．
補題~\ref{lem:group-heavy-antichain} より
重い葉は高々 $M_{A,f}-1$ 個であり，
各重い節点はそのいずれかへの道上にある．
したがって重い節点の総数は高々
$(M_{A,f}-1)n_g$ である．
根が軽ければ長さはすでに $\bar\theta_G$ 以下である．
そうでなければ，各重い節点の規則が明示的に出力する
終端数は高々 $\bar\theta_G$，
その高々 $r_{\max}$ 個の軽い子部分木の総長も
各々高々 $\bar\theta_G$ である．
重い節点と極大軽部分木を数えると結論が従う．
\end{proof}

\begin{lemma}[窓と任意有限群の型付けに対する短い正準タプル]
\label{lem:window-group-short}
定理~\ref{thm:group-typed-shortest} の仮定のもと，
$k,\ell\ge0$ と
$h\preceq(h_\Sigma^{k,\ell},g)$ を満たす
固定有限型付け $h$ を選ぶ．
$G$ の刈り込み済み完全 $h$-精密化を
$\widetilde G_0$ とし，その開始記号以外の
非終端記号数を $n_{\mathrm{typ}}$ とする．
各非開始記号 $X$ について，
\[
 \|\omega(X)\|_1\le
 f(k+\ell+1)n_{\mathrm{typ}}
 \bigl(\bar\theta_G+r_{\max}\beta_g^{\mathrm{grp}}\bigr)
\]
を満たす正準タプルを選べる．
\end{lemma}
\begin{proof}
$X$ からの終端導出で節点数が最小のものを選ぶ．
根から葉への道上に同じ型付き非終端記号が
二度現れれば，根の $h$-型を保存したまま
節点数を減らせるため，各道は
$n_{\mathrm{typ}}$ 節点以下である．
補題~\ref{lem:boundary-short-canonical} と同様に，
各根成分で $h_\Sigma^{k,\ell}$ の値を決める
位置に印を付ける．短い成分では全位置，
長い成分では先頭 $k$ 位置と末尾 $\ell$ 位置を，
$k=\ell=0$ なら一点を保護する．
高々 $f(k+\ell+1)$ 個の印への祖先経路が作る
骨格の節点数は高々
$f(k+\ell+1)n_{\mathrm{typ}}$ である．

骨格から外れる各子部分木の，
元の非終端記号 $B$ と実際の成分別群値を記録する．
定理~\ref{thm:group-typed-shortest} により，
同じ $B$ から，同じ各成分群値を持ち，
総長が $\beta_g^{\mathrm{grp}}$ 以下の代替導出を選べる．
$G$ は良形なので代替成分はすべて非空である．
保護位置は変化せず，長い根成分は長いままであるため，
窓の値は保存される．また，各子成分の群値を
個別に保存しているので，規則テンプレートを
元の順序で評価すれば根の各群値も保存される．
可換性は必要ない．
したがって積型付け $(h_\Sigma^{k,\ell},g)$，
およびより粗い根の $h$-型は不変である．
命題~\ref{prop:output-invariants}(ii) により，
元の $G$-導出を，根を $X$ とする
刈り込み済み完全精密化の導出へ持ち上げられる．

各骨格節点が明示する終端は高々
$\bar\theta_G$ 個であり，
高々 $r_{\max}$ 個の省略子それぞれは
長さ高々 $\beta_g^{\mathrm{grp}}$ を寄与する．
ゆえに上界が従う．
\end{proof}

\begin{theorem}[窓と任意有限群による型付けの多項式特徴データ]
\label{thm:window-group-polydata}
$f,r_{\max},k,\ell,\Sigma$，任意の有限群 $A$，
準同型 $g:\Sigma^*\to A$ および
$h\preceq(h_\Sigma^{k,\ell},g)$ を満たす固定有限型付け
$h$ を固定する．
$\Cmcf{f}{h}$ の言語を生成する任意の
reduced good fan-out-$\le f$，
rank-$\le r_{\max}$ 文法 $G$ に対して，
ランク制限つき再構成作用素
$\mathcal B^{\le r_{\max}}_{f,h}$ は
次を満たす特徴標本を持つ：
\[
 |\Wit(\widetilde G_0)|=O(\|G\|),\qquad
 \|\Wit(\widetilde G_0)\|_+
 =O_{f,r_{\max},k,\ell,h,g,\Sigma}
   (\|G\|^3\bar\theta_G).
\]
\end{theorem}
\begin{proof}
補題~\ref{lem:window-group-short} より，
刈り込み済み完全精密化の正準タプルについて
\[
 b_{\mathrm{can}}\le
 f(k+\ell+1)n_{\mathrm{typ}}
 (\bar\theta_G+r_{\max}\beta_g^{\mathrm{grp}})
 =O(\|G\|^2\bar\theta_G)
\]
を得る．ここで型付けと規則ランクの上界は固定され，
$n_{\mathrm{typ}}\le|V||M|^f=O(\|G\|)$ である．
補題~\ref{lem:untyped-contexts} より
各特徴証拠語の長さは
$2\bar\theta_G+r_{\max}b_{\mathrm{can}}
=O(\|G\|^2\bar\theta_G)$ 以下となり，
命題~\ref{prop:witness-cardinality} より
語数は $O(\|G\|)$ 以下である．
空語の証拠を加えても総符号化長は定数しか増えない．
\end{proof}

\begin{remark}
この強化は，任意の有限\emph{モノイド}型付けについて
元文法の厚さから多項式上界が得られると主張しない．
$f$ を固定すると，定数
$M_{A,f}=U_f\max\{3,|A|^f\}^{2^{f^2}}$ は $|A|$ について多項式である．
これは以前の版で用いた多色Ramsey上界
$U_f\,c_f^{\,c_f(N_{A,f}-1)+1}$（$c_f=2^{f^2}$）を置き換える．
後者は $|A|^f$ について指数的であった．
ただし $f$ への依存は二重指数的なままである．
$A$ が可換な場合は，\ref{app:window-abelian} の直接証明が
$f$ について線形な，はるかに小さい定数を与える
（注意~\ref{rem:davenport-size}）．
\end{remark}
